% GENERATED FILE. Edit canonical lessons, then run npm run generate:books.
% canonical-source-hash: fnv1a64:84e7c525fb826bbb
% canonical-lessons: RU-C515-royal, RU-C515-violonchel, RU-C515-fleyta, RU-C515-struna, RU-C515-nastraivat

\chapter{Piano, Cello and Flute}
\label{ch:ru-piano-cello-and-flute}
\hlchaptermodality{515}

\begin{chapteropening}
\textbf{By the end of this chapter:} \emph{I can talk about the instruments people play.}

Complete the last lesson of chapter 515: I can talk about the instruments people play.
\end{chapteropening}

\section[\texorpdfstring{\ru{рояль} (royál)}{рояль (royál)}]{\ru{рояль} (royál) — a grand piano}

\label{lesson:RU-C515-royal}



\begin{quote}
Before the new one: say the Russian for flippers, then the Russian for to splash about.
\end{quote}

\subsection*{You'll want to know: \ru{рояль}}
\textbf{\ru{рояль}} (\emph{royál}) — \textquotedblleft{}a grand piano\textquotedblright{}.

\begin{quote}
\ru{В} \ru{зале} \ru{стоит} \ru{большой} \ru{чёрный} \ru{рояль}.
\end{quote}

\begin{quote}
\emph{V zále stoít bolshóy chyórnyy royál.}
\end{quote}

\begin{quote}
There is a big black grand piano in the hall.
\end{quote}

Masculine.

\begin{cousinweb}
EtymDB lists \emph{\ru{рояль}}, glossed 'grand piano', without a source; it is usually traced to French \emph{royal} 'royal'.
\end{cousinweb}

\begin{grammarlens}[title={What you've built}]
One more word for telling, explaining and describing at B1.
\end{grammarlens}

\subsection*{Your turn}
\begin{itemize}
\raggedright
  \item \emph{Say it:} \emph{royál}
  \item \emph{Say it:} \emph{royál}, once more
  \item \emph{Say it:} a sentence of your own with \emph{royál}
  \item \emph{Recall:} say the Russian for flippers, then the Russian for to splash about, then say \emph{royál} again
\end{itemize}

\begin{tcolorbox}[breakable,colback=teal!4,colframe=teal!35!black,title={Before you move on}]
What does \ru{рояль} mean? (\textquotedblleft{}A grand piano\textquotedblright{}.) Say it once more.
\end{tcolorbox}
\section[\texorpdfstring{\ru{виолончель} (violonchél)}{виолончель (violonchél)}]{\ru{виолончель} (violonchél) — a cello}

\label{lesson:RU-C515-violonchel}



\begin{quote}
Before the new one: say the Russian for to splash about, then the Russian for a grand piano.
\end{quote}

\subsection*{You'll want to know: \ru{виолончель}}
\textbf{\ru{виолончель}} (\emph{violonchél}) — \textquotedblleft{}a cello\textquotedblright{}.

\begin{quote}
\ru{Моя} \ru{сестра} \ru{играет} \ru{на} \ru{виолончели}.
\end{quote}

\begin{quote}
\emph{Moyá sestrá igráyet na violonchéli.}
\end{quote}

\begin{quote}
My sister plays the cello.
\end{quote}

Feminine. \emph{\ru{Играть} \ru{на} \ru{виолончели}}.

\begin{cousinweb}
\emph{\ru{виолончель}} comes from Italian \emph{violoncello} (EtymDB), a diminutive of \emph{violone} 'double bass'.
\end{cousinweb}

\begin{grammarlens}[title={What you've built}]
One more word for telling, explaining and describing at B1.
\end{grammarlens}

\subsection*{Your turn}
\begin{itemize}
\raggedright
  \item \emph{Say it:} \emph{violonchél}
  \item \emph{Say it:} \emph{violonchél}, once more
  \item \emph{Say it:} a sentence of your own with \emph{violonchél}
  \item \emph{Recall:} say the Russian for to splash about, then the Russian for a grand piano, then say \emph{violonchél} again
\end{itemize}

\begin{tcolorbox}[breakable,colback=teal!4,colframe=teal!35!black,title={Before you move on}]
What does \ru{виолончель} mean? (\textquotedblleft{}A cello\textquotedblright{}.) Say it once more.
\end{tcolorbox}
\section[\texorpdfstring{\ru{флейта} (fléyta)}{флейта (fléyta)}]{\ru{флейта} (fléyta) — a flute}

\label{lesson:RU-C515-fleyta}



\begin{quote}
Before the new one: say the Russian for a grand piano, then the Russian for a cello.
\end{quote}

\subsection*{You'll want to know: \ru{флейта}}
\textbf{\ru{флейта}} (\emph{fléyta}) — \textquotedblleft{}a flute\textquotedblright{}.

\begin{quote}
\ru{Звук} \ru{флейты} \ru{очень} \ru{нежный}.
\end{quote}

\begin{quote}
\emph{Zvuk fléyty óchen nézhnyy.}
\end{quote}

\begin{quote}
The sound of the flute is very gentle.
\end{quote}

Feminine.

\begin{cousinweb}
\emph{\ru{флейта}} was borrowed from German \emph{Flöte} (EtymDB).
\end{cousinweb}

\begin{grammarlens}[title={What you've built}]
One more word for telling, explaining and describing at B1.
\end{grammarlens}

\subsection*{Your turn}
\begin{itemize}
\raggedright
  \item \emph{Say it:} \emph{fléyta}
  \item \emph{Say it:} \emph{fléyta}, once more
  \item \emph{Say it:} a sentence of your own with \emph{fléyta}
  \item \emph{Recall:} say the Russian for a grand piano, then the Russian for a cello, then say \emph{fléyta} again
\end{itemize}

\begin{tcolorbox}[breakable,colback=teal!4,colframe=teal!35!black,title={Before you move on}]
What does \ru{флейта} mean? (\textquotedblleft{}A flute\textquotedblright{}.) Say it once more.
\end{tcolorbox}
\section[\texorpdfstring{\ru{струна} (struná)}{струна (struná)}]{\ru{струна} (struná) — a string (of an instrument)}

\label{lesson:RU-C515-struna}



\begin{quote}
Before the new one: say the Russian for a cello, then the Russian for a flute.
\end{quote}

\subsection*{You'll want to know: \ru{струна}}
\textbf{\ru{струна}} (\emph{struná}) — \textquotedblleft{}a string (of an instrument)\textquotedblright{}.

\begin{quote}
\ru{На} \ru{гитаре} \ru{порвалась} \ru{струна}.
\end{quote}

\begin{quote}
\emph{Na gitáre porvalás struná.}
\end{quote}

\begin{quote}
A string has broken on the guitar.
\end{quote}

Feminine; plural \ru{струны}.

\begin{cousinweb}
\emph{\ru{струна}} is inherited from Proto-Slavic \emph{struna} (EtymDB).
\end{cousinweb}

\begin{grammarlens}[title={What you've built}]
One more word for telling, explaining and describing at B1.
\end{grammarlens}

\subsection*{Your turn}
\begin{itemize}
\raggedright
  \item \emph{Say it:} \emph{struná}
  \item \emph{Say it:} \emph{struná}, once more
  \item \emph{Say it:} a sentence of your own with \emph{struná}
  \item \emph{Recall:} say the Russian for a cello, then the Russian for a flute, then say \emph{struná} again
\end{itemize}

\begin{tcolorbox}[breakable,colback=teal!4,colframe=teal!35!black,title={Before you move on}]
What does \ru{струна} mean? (\textquotedblleft{}A string (of an instrument)\textquotedblright{}.) Say it once more.
\end{tcolorbox}
\section[\texorpdfstring{\ru{настраивать} (nastráivat)}{настраивать (nastráivat)}]{\ru{настраивать} (nastráivat) — to tune; to set up}

\label{lesson:RU-C515-nastraivat}



\begin{quote}
Before the new one: say the Russian for a flute, then the Russian for a string.
\end{quote}

\subsection*{You'll want to know: \ru{настраивать}}
\textbf{\ru{настраивать}} (\emph{nastráivat}) — \textquotedblleft{}to tune; to set up\textquotedblright{}.

\begin{quote}
\ru{Перед} \ru{концертом} \ru{скрипач} \ru{настраивает} \ru{скрипку}.
\end{quote}

\begin{quote}
\emph{Péred kontsértom skripách nastráivayet skrípku.}
\end{quote}

\begin{quote}
Before the concert the violinist tunes his violin.
\end{quote}

Imperfective; perfective \ru{настроить}. Also to set up a device.

\begin{cousinweb}
\emph{\ru{настраивать}} is built on \emph{\ru{на}-} and \emph{\ru{строить}} 'to build, to arrange'; EtymDB lists the verb without a source.
\end{cousinweb}

\begin{grammarlens}[title={What you've built}]
One more word for telling, explaining and describing at B1.
\end{grammarlens}

\subsection*{Your turn}
\begin{itemize}
\raggedright
  \item \emph{Say it:} \emph{nastráivat}
  \item \emph{Say it:} \emph{nastráivat}, once more
  \item \emph{Say it:} a sentence of your own with \emph{nastráivat}
  \item \emph{Recall:} say the Russian for a flute, then the Russian for a string, then say \emph{nastráivat} again
\end{itemize}

\begin{tcolorbox}[breakable,colback=teal!4,colframe=teal!35!black,title={Before you move on}]
What does \ru{настраивать} mean? (\textquotedblleft{}To tune; to set up\textquotedblright{}.) Say it once more.
\end{tcolorbox}
